\chapter{An event begins with its baseline}\label{ch:event}
The same physical increment can receive different values at different
states. It follows that the event boundary is part of the scientific
definition, not merely a scheduler detail. A record must identify when
forcing occurs, when the state is frozen and which version each candidate
uses.

\section{The intended forcing-first order}
The new programme proposes the following logical sequence:
\begin{center}
\begin{tabular}{P{.13\linewidth}P{.74\linewidth}}
\toprule Phase & Responsibility\\\midrule
1 & Apply the declared external forcing and record its signed audit.\\
2 & Freeze one post-forcing physical baseline and its configuration.\\
3 & Generate EBU-blind candidates/groups and check joint physical feasibility.\\
4 & Evaluate exact group value and declared path attribution.\\
5 & Aggregate by declared owner and screen affordability.\\
6 & Make an unweighted random choice from the declared eligible menu.\\
7 & Execute exactly that selected accepted vector, then record accounting effects once.\\
\bottomrule
\end{tabular}
\end{center}
These are explanatory phases, not a new accepted enumeration of framework
phase constants. The inspected framework has its own fixed ten-phase
chronology, including natural drive in a different position. A recognized
Gaussian event profile must resolve that difference before runtime adoption
\cite{framework}.

\section{Why forcing is separate}
Write $x_t$ for the pre-forcing state and $z_t$ for the post-forcing baseline.
The external record is $D_{\mathrm{ext},t}=V(z_t)-V(x_t)$. All candidates
in that epoch compare $z_t$ with their own declared endpoints. An idle actor
therefore has zero action value even if forcing just improved the field.

If one candidate uses $x_t$ and another uses $z_t$, their comparison includes
different external effects. If forcing changes again between quote and
execution, the committed conditions need explicit treatment; the software
cannot quietly call the new state the old baseline.

\section{Epoch identity binds meaning}
A useful epoch commitment identifies the state version, physical profile,
potential and parameter version, accepted child vector, complete support,
path semantics, owner-rule version and numerical policy. Not every module
needs to read all of these values, but the composite record must connect
them unambiguously.

Changing an accepted quantity changes the endpoint and generally every
shared-path attribution. Removing one child can change the value of the
others. Such changes require a new affected valuation and authorization
record, not reuse of an old number whose identity happens to look familiar.

\section{A worked baseline error}
At reference $(10,10)$ the potential is zero. Forcing produces $(14,6)$
with potential four. A restoring action ends at $(12,8)$ with potential one.
The actor's value is $4-1=3$. Comparing the original reference directly
with the endpoint gives $0-1=-1$, which includes the shock and answers a
different question. Both subtractions are arithmetically correct; only one
is the declared actor account.

This is why a test of subtraction alone is insufficient. The implementation
must also test which state commitment supplies each operand.

\section{Failure and publication order}
A prospective runtime must specify what is durable before state advancement
becomes possible, what constitutes one completed event and how partial
publication is detected. Book prose cannot substitute for that execution
contract. Historical Gate 1D-C has a particular durability protocol; it
must not be copied as though every new study shares its exact permissions.

\section*{Chapter recap}
The baseline, chronology and event identity are part of the model. A new
diagram explains the intended process but does not amend the framework's
existing scientific order by itself.
